\documentclass[10pt]{article}
\usepackage[margin=12mm]{geometry}
\usepackage{graphicx}
\usepackage{caption}
\pagestyle{empty}
\captionsetup{font=small,labelfont=bf}

\begin{document}
\begin{figure}[p]
\centering
\begin{minipage}[t]{0.485\textwidth}
  \centering
  \includegraphics[width=\linewidth]{figS1A_ce_cousin_r_definition.pdf}
\end{minipage}\hfill
\begin{minipage}[t]{0.485\textwidth}
  \centering
  \includegraphics[width=\linewidth]{figS1B_ce_cousin_r_major_subtrees.pdf}
\end{minipage}

\vspace{2mm}
\includegraphics[width=\textwidth]{figS1C_ce_cousin_r_all_subtrees.pdf}

\caption{\textbf{Sensitivity of subtree proximity to the first-cousin null.}
%
\textbf{(A)} The first-cousin-null-relative construction. $P^*$ is the sampled
Pareto assignment closest to the natural lineage $L$ in separately
first-cousin-null-standardized travel and cell-state coordinates, and
$r=\left\|L-P^*\right\|_2/\left\|N-P^*\right\|_2$, where $N$ is the
first-cousin-null mean. Thus $r=0$ means lineage is on the Pareto front and $r=1$ means it is at the null
reference.
%
\textbf{(B)} Canonical positions of the major subtrees P0, AB, ABa, ABp, and P1 using
the cousin-null-relative metric.
%
\textbf{(C)} Cousin-null-relative positions of all eligible subtrees. Purple
markers show the natural-lineage metric, and shape follows the lineage-region
encoding in Figure 5. The
inset compares $r$ with the null-independent endpoint-normalized distance
$d_{LP}$; their strong rank agreement (Spearman $\rho=0.95$) shows that the
broad natural-lineage-to-Pareto-front proximity ranking is robust regardless of the
null reference.}
\label{fig:ce_canonical_summary_cousin_r}
\end{figure}
\end{document}
